\documentclass[11pt,a4paper]{article}
% Common preamble for the hint sheets (\input from each file)
\usepackage[utf8]{inputenc}
\usepackage[T1]{fontenc}
\usepackage{lmodern}
\usepackage{amsmath,amssymb,amsthm,mathtools}
\usepackage[a4paper,margin=2.3cm]{geometry}
\usepackage{enumitem}
\usepackage{xcolor}
\usepackage[hidelinks]{hyperref}
\renewcommand{\leq}{\leqslant}
\renewcommand{\le}{\leqslant}
\renewcommand{\geq}{\geqslant}
\renewcommand{\ge}{\geqslant}
\setlength{\parindent}{0pt}
\setlength{\parskip}{0.35em}
% \hintstitle{Sheet title}{Class work / Homework}
\newcommand{\hintstitle}[2]{\begin{center}{\LARGE\bfseries Hints: #1\par}\vspace{0.5em}{\large #2\par}\end{center}
\noindent\textit{The hints to each problem are cumulative. Read one, try again, and open the next only if you are still stuck.}\par\medskip\hrule\medskip}
\newcommand{\problem}[1]{\par\bigskip{\large\bfseries Problem #1}\par\nopagebreak\smallskip}
\newcommand{\subpart}[1]{\par\smallskip\textbf{(#1)}\par\nopagebreak}
\newcommand{\hint}[1]{\par\smallskip\noindent\textcolor{black!70}{\textbf{Hint #1.}}\enspace}
\newcommand{\note}[1]{\par\smallskip\noindent\textit{Note.} #1\par}
\providecommand{\lt}{<}
\providecommand{\gt}{>}
\providecommand{\tightlist}{\setlength{\itemsep}{0pt}\setlength{\parskip}{0pt}}
\emergencystretch=3em

\begin{document}
\hintstitle{Linear algebra -- 9. Matrices over \(\mathbb Z\) and over a finite field}{Homework}

\problem{1}
\hint{1}
An invertible matrix is a list of columns, each outside the span of the previous ones: there are \(q^n-q^k\) choices for column \(k+1\). For subspaces, count ordered \(k\)-tuples of independent vectors in two ways: directly, and as (number of subspaces)\(\times|\mathrm{GL}_k(\mathbb F_q)|\). This gives \(\prod_{i<k}\frac{q^n-q^i}{q^k-q^i}\).
\hint{2}
\(X=|\ker(g-E)|\), and \(X^k\) is the number of \(k\)-tuples \((v_1,\dots,v_k)\) fixed by \(g\). By Burnside's lemma \(\mathsf EX^k\) is the number of orbits of \(\mathrm{GL}_n(\mathbb F_q)\) on \((\mathbb F_q^n)^k\).
\hint{3}
A \(k\)-tuple is a linear map \(\mathbb F_q^k\to\mathbb F_q^n\). For \(n\ge k\) two tuples lie in one orbit if and only if they have the same kernel (extend a basis of the image to a basis of \(\mathbb F_q^n\)). So the orbits correspond to the subspaces of \(\mathbb F_q^k\), and \(\mathsf EX^k=\sum_{j=0}^k\binom kj_q\). For \(k=1\): \(\mathsf EX=2\); for \(k=2\): \(\mathsf EX^2=1+(q+1)+1=q+3\), so \(\operatorname{Var}X=q-1\).

\problem{2}
\hint{1}
The polynomial \(x^n+1\) has distinct roots and annihilates \(A\). So \(A\) is diagonalisable over \(\mathbb C\), and its eigenvalues are roots of unity with \(\lambda^n=-1\). The characteristic polynomial is a rational quadratic.
\hint{2}
If the eigenvalues are rational, both equal \(-1\) and \(A=-E\). Otherwise the characteristic polynomial is a cyclotomic polynomial of degree \(2\): \(\Phi_3\), \(\Phi_4\) or \(\Phi_6\). Rule out \(\Phi_3\) (no cube root of unity has a power equal to \(-1\)). By Cayley--Hamilton, \(\Phi_4\) gives \(A^2=-E\) and \(\Phi_6\) gives \(A^3=-E\).

\problem{3}
\hint{1}
For \(X\in\mathrm{SL}_2(\mathbb Z)\), Cayley--Hamilton gives \(X^2=(\operatorname{tr}X)X-E\). The traces drive everything.
\subpart{a}
\hint{1}
Try \(\operatorname{tr}A=-1\), \(\operatorname{tr}B=1\), \(\operatorname{tr}C=0\). Then \(A^2+B^2=B-A-2E\) and \(C^2=-E\), so you need \(B=A+E\). For example \(A=\begin{pmatrix}-1&-1\\1&0\end{pmatrix}\), \(B=\begin{pmatrix}0&-1\\1&1\end{pmatrix}\), \(C=\begin{pmatrix}-1&-2\\1&1\end{pmatrix}\).
\subpart{b}
\hint{1}
Take traces: \(\operatorname{tr}X^4=t^4-4t^2+2\) for \(t=\operatorname{tr}X\). This gives \(a^4+b^4-c^4-4(a^2+b^2-c^2)+2=0\) for the traces \(a,b,c\).
\hint{2}
Modulo \(4\), fourth powers are \(0\) or \(1\), so \(a,b\) are odd and \(c\) is even. Then work modulo \(16\): \(a^2+b^2\equiv2\pmod 8\) and \(c^2\equiv0\) or \(4\pmod 8\). The left side is \(\equiv4-8\not\equiv0\pmod{16}\).

\problem{4}
\hint{1}
The matrix is \(J+D\) with \(D=\operatorname{diag}(k^2-1)_{k=2}^{N}\), \(N=p+7\). So \(\det=\prod(k^2-1)\cdot\bigl(1+\sum 1/(k^2-1)\bigr)\). Telescope both.
\hint{2}
You should get \(\det=\bigl((N-1)!\bigr)^2\,\dfrac{7N^2+3N-2}{8}\). For \(p\ge7\), \(p\) divides \((p+6)!\) exactly once. So you need \(p\mid 7\cdot7^2+3\cdot7-2=362=2\cdot181\). Check \(p=2,3,5\) directly.
\hint{3}
Answer: \(p\in\{2,3,5,181\}\).

\problem{5}
\hint{1}
The entries \(a_{\sigma(i)i}\) lie in distinct rows and columns. Hence \(f(\varepsilon)=\det\) of the modified matrix is a polynomial in \(\varepsilon(1),\dots,\varepsilon(n)\) of degree at most \(1\) in each variable. What is the coefficient of \(\varepsilon(1)\cdots\varepsilon(n)\)?
\hint{2}
That coefficient is \(\operatorname{sgn}\sigma=\pm1\). For a multilinear polynomial, \(\sum_{\varepsilon\in\{0,1\}^n}(-1)^{n-|\varepsilon|}f(\varepsilon)\) equals this top coefficient. If \(m\) divided every \(f(\varepsilon)\), it would divide \(\pm1\).

\problem{6}
\hint{1}
For \(p=2\) compare orders. For odd \(p\), look at \(u=\begin{pmatrix}1&1\\0&1\end{pmatrix}\) of order \(p\). In \(S_p\) an element of order \(p\) is a \(p\)-cycle. What is its centraliser?
\hint{2}
The centraliser of a \(p\)-cycle in \(S_p\) has order \(p\). The centraliser of \(u\) in \(\mathrm{GL}_2(\mathbb Z/p)\) contains all \(\begin{pmatrix}a&b\\0&a\end{pmatrix}\), which is \(p(p-1)\) elements. An injective homomorphism maps \(C(u)\) injectively into \(C(\varphi(u))\).

\problem{7}
\hint{1}
Take logarithms of the (positive) entries. The operations become: add the other row (column) to a row (column), or subtract it.
\hint{2}
These operations preserve the determinant of the matrix of logarithms. It is \(\ln2(\ln4-\ln3)\) at the start and its negative at the target, and it is non-zero because \(\ln4\ne\ln3\). Answer: no.

\problem{8}
\hint{1}
Rank-one reduction (Wedderburn): \(\operatorname{rank}(A-ae_ie_j^T)<\operatorname{rank}A\) if and only if \(e_i\in\operatorname{Im}A\), \(e_j^T\) is in the row space and \(a\,e_j^TGe_i=1\), where \(G\) is any generalised inverse (\(AGA=A\)). For invertible \(A\) this says \((A^{-1})_{ji}=1/a_{ij}\).
\hint{2}
Read off the \((i,i)\) entry of the idempotent \(AG\). Since \(e_i=Ax_0\), it equals \(1\); on the other hand it is \(\sum_j a_{ij}G_{ji}\), a sum of ones over the non-zero entries of row \(i\). So the number of non-zero entries of each non-zero row is \(\equiv1\pmod p\). For \(p>2014\) this forces exactly one.
\hint{3}
For \(p\le2013\) a counterexample: \(J-E\) of size \(m=p+2\) has inverse \(\frac{1}{m-1}J-E=J-E\) over \(\mathbb F_p\). Put it in a block-diagonal matrix with an identity block. Answer: exactly the primes \(p\ge2017\).

\problem{9}
\hint{1}
For \(j\ge1\), \(i^j\bmod3\) depends only on \(i\bmod3\) and the parity of \(j\). Show that rows \(i\) and \(i+6\) coincide, so \(\det A=0\) for \(n\ge7\).
\hint{2}
Compute the small cases by hand (only residues \(0,1,2\) appear) and find a dependency among the rows for \(n=6\). Answer: \(n=5\) (indeed \(\det A=12\) there).

\problem{10}
\hint{1}
Orient the edges and let \(D\) be the \(n\times m\) incidence matrix, so that \(L=DD^T\). Delete a row of \(D\) to get \(D_0\) and apply Cauchy--Binet to \(\det(D_0D_0^T)\).
\hint{2}
An \((n-1)\times(n-1)\) minor of \(D_0\) is \(\pm1\) if the chosen edges form a spanning tree, and \(0\) otherwise. A cycle gives dependent columns; for a tree, expand at a leaf.
\hint{3}
For \(L(K_n)=nE-J\): the entries generate \(1\), so \(d_1=1\). Modulo \(n\), \(L\equiv-J\) has rank \(1\), so all \(k\times k\) minors are divisible by \(n^{k-1}\). Moreover \(d_1\cdots d_{n-1}=n^{n-2}\) (Cayley). Hence the Smith form is \(\operatorname{diag}(1,n,\dots,n,0)\) and the critical group is \((\mathbb Z/n)^{n-2}\).
\end{document}
